\section{Introduction}

DBpedia extracts structured information from Wikipedia---infobox fields,
abstracts, categories, redirects and links between articles---and publishes it
as RDF, so that the content of an encyclopedia can be queried with SPARQL and
joined with other datasets~\cite{auer2007dbpedia, lehmann2015dbpedia}. The
English edition sits at the centre of the Linked Open Data cloud, and language
chapters publish the same kind of data for other Wikipedia
editions~\cite{kontokostas2012greek}. Vietnamese Wikipedia articles carry the
same semi-structured information: the infobox of a footballer lists his date of
birth, birthplace and every club he played for, season by season. As wikitext,
however, it cannot be queried. A question such as \emph{``Which players born in
Nghệ An have played for the national team, and for which club do they play
now?''} requires reading and joining several articles by hand.

This project addresses Topic~2 of the course, \emph{Build a DBpedia version for
Vietnamese language}, which consists of five tasks: (1)~define an ontology;
(2)~collect Vietnamese articles from Wikipedia; (3)~transform the data into
4-star Linked Data; (4)~establish links to the English DBpedia; and
(5)~provide an interface to query the data through a SPARQL endpoint or
terminal.

Rather than extracting all of Vietnamese Wikipedia shallowly, we chose a closed
domain and extracted it in depth: Vietnamese \emph{football} (players, clubs,
national teams, stadiums), \emph{provinces} and \emph{universities}. Three
properties of this domain motivated the choice. First, its entities refer to
one another, so the graph supports multi-hop questions: player $\rightarrow$
career station $\rightarrow$ club $\rightarrow$ home stadium $\rightarrow$
province, player $\rightarrow$ birth province, university $\rightarrow$
province. Second, 777 of the 990 selected entities (78\%) have an English
Wikipedia article, so most of them can be linked to English DBpedia. Third,
956 of the 990 articles (97\%) contain an infobox, the main source that DBpedia
extracts. Entities are selected through Wikidata rather than through the
category tree of Vietnamese Wikipedia, because categories are assigned by
editors and are not used consistently.

The contributions of this work are:
\begin{itemize}
  \item an OWL ontology, \term{vio:}, with 18 classes and 54 properties. Every
        class has a DBpedia Ontology ancestor, 38 properties are
        sub-properties of a \term{dbo:} property, and OWL~2~RL axioms
        (a property chain, inverses, restrictions and disjointness) are
        designed for reasoning (Section~\ref{sec:ontology});
  \item a reproducible extraction pipeline of six commands. All network
        requests go through a local cache, and the raw downloads are
        committed, so the dataset can be rebuilt without Internet access
        (Sections~\ref{sec:collection}--\ref{sec:postprocess});
  \item a dataset of 131{,}233 triples, of which 88{,}898 are asserted and
        41{,}730 are inferred, that passes eight automated quality checks with
        no error or warning (Section~\ref{sec:evaluation});
  \item 777 \term{owl:sameAs} links to English DBpedia, verified against the
        live endpoint, and 990 to Wikidata (Sections~\ref{sec:interlinking}
        and~\ref{sec:links});
  \item a server that dereferences every resource IRI with HTTP 303 and
        content negotiation, and a web interface with resource pages, a SPARQL
        editor and natural-language question answering
        (Section~\ref{sec:publishing}).
\end{itemize}

Table~\ref{tab:requirements} maps each task of the assignment to the part of
the system that implements it. Section~\ref{sec:related} introduces DBpedia
and the Linked Data principles we follow. Section~\ref{sec:design} describes
the ontology and the extraction pipeline, Section~\ref{sec:publishing} the
server and the query interfaces, and Section~\ref{sec:evaluation} evaluates the
dataset. Section~\ref{sec:discussion} discusses design decisions and
limitations, and Section~\ref{sec:conclusion} concludes.

\begin{table}[t]
\centering
\caption{Tasks of the assignment and where they are implemented. Module paths
are relative to \term{vidbpedia/}.}
\label{tab:requirements}
\small
\begin{tabularx}{\textwidth}{@{}l>{\raggedright\arraybackslash}X>{\raggedright\arraybackslash}Xl@{}}
\toprule
Task & Implementation & Output & Section \\
\midrule
1. Ontology & \term{ontology/*.ttl}, merged by \term{kg/ontology.py}
  & 18 classes, 54 properties, 549 triples & \ref{sec:ontology} \\
2. Collect articles & \term{crawl/wikidata\_seeds.py}, \term{crawl/wiki\_enrich.py}
  & 990 entities with their Vietnamese articles & \ref{sec:collection} \\
3. 4-star data & \term{crawl/build\_rdf.py}, \term{kg/postprocess.py}
  & 131{,}233 triples in Turtle and N-Triples & \ref{sec:transform}, \ref{sec:postprocess} \\
4. Link to DBpedia & enwiki sitelinks of the Wikidata item
  & 777 \term{owl:sameAs dbr:} links, VoID linksets & \ref{sec:interlinking} \\
5. Query interface & \term{web/}: FastAPI and Gradio
  & Linked Data routes, SPARQL editor, question answering & \ref{sec:publishing} \\
\bottomrule
\end{tabularx}
\end{table}
